\documentstyle[%


righttag%


,11pt%


,amssymb%


,russian%               remove in Eng.


,izvru%                 Set izven% in Eng.


]{amsart}





% 


\newcommand{\tld}{\widetilde}


\renewcommand{\phi}{\varphi}


\newcommand{\norm}[1]{\left\|#1\right\|}


\newcommand{\abs}[1]{\left|#1\right|}


\newcommand{\set}[1]{\left\{#1\right\}}


\newcommand{\supp}{\operatorname{supp}}


\newcommand{\Lone}{L_{1}}


\newcommand{\Ltwo}{L_{2}}


\newcommand{\tts}[1]{}


\numberwithin{equation}{section}





\theoremstyle{plain}


\theorembodyfont{\it}


\newtheorem{thms}{\hskip\parindent Теорема}[section]


\newtheorem{lems}{\hskip\parindent Лемма}[section]


%% \renewcommand{\thethmnonum}{}





\theoremstyle{definition}


\theorembodyfont{\rm}


\newtheorem{defnnonum}{\hskip\parindent Определение}


\renewcommand{\thedefnnonum}{}





%\hoffset=-15mm%                  For Eng.


\setcounter{page}{52}


\god{2005}


\nomer{9~(520)} %                        Remove in Eng.


%\nomer{Vol.\,49, No.\,8}%               Set in Eng.


%\str{\pageref{begin}--\pageref{end}}%  Set in Eng.


\udk{517.983}





\author{\it А.Ю.\,Новосельцев, И.Б.\,Симоненко}


% NB:   ^^ remove in Eng.


\title{Зависимость асимптотики старших собственных значений


усеченной континуальной свертки от скорости достижения максимума символом}





\date{}


%\pagestyle{myheadings}%         Set in Eng.


\pagestyle{plain} %              Remove in Eng.


\begin{document}


%\thispagestyle{plain}%          Set in Eng.


\maketitle


\label{begin}





\setcounter{section}{0}





В работе исследуется асимптотическое поведение крайних собственных 


значений оператора усеченной континуальной свертки


\begin{equation} \label{operatorAy}


[A_{y} f](t)=\int_{0}^{y}k(t-s)f(s)ds,\ \ y>0,\ \ t\in[0,y],


\end{equation}


с вещественнозначным символом $K(x)$, достигающим своего максимума 


в одной или в конечном числе точек.





Основными результатами данной работы являются теоремы \ref{th:main}


и \ref{th:general}, в которых установлено, что старшие собственные 


числа оператора $A_{y}$ при $y\to+\infty$ стремятся к


$M=\max\limits_{x\in\Bbb R}K(x)$ со скоростью ${1}/{y^{\nu}}$, где


$\nu>0$ --- порядок нуля функции $M-K(x)$. 





Исследование представляет собой как бы ответ на работу \cite{li:Serra},


где аналогичный результат для оператора усеченной дискретной свертки был 


получен лишь для четных $\nu$, а в других случаях доказаны


оценки, использующие окаймляющие $\nu$ четные числа. При этом автор 


опирался на основополагающие результаты работ


\cite{li:Szego1920}--\cite{li:Szego1952}, содержащиеся также в 


\cite{li:SzegoForms} (русский перевод предыдущего издания 


этой книги --- \cite{li:SzegoRus}). Наш метод не использует их и 


позволяет получить точные по порядку оценки сверху и снизу для любых $\nu>0$. 








\section{Обозначения}





Будем пользоваться следующими обозначениями: $\Bbb N$ --- множество 


натуральных чисел; $\Bbb R$, $\Bbb C$ --- поля вещественных и комплексных 


чисел соответственно; $\bold i$ --- мнимая единица, $\bold e$ --- основание 


натурального логарифма; $\mu$ --- мера Лебега на $\Bbb R$.





Через $\cal F$ обозначим преобразование Фурье, действующее на функцию 


$f\in\Ltwo(\Bbb R)$ следующим образом:


$[\cal F f](x)=\frac{1}{\sqrt{2\pi}}\int\limits_{\Bbb R}f(t)


\bold e^{-\bold i xt}dt$.


Хорошо известно, что преобразование Фурье обратимо и


$\norm{f}_{\Ltwo}=\norm{\cal F f}_{\Ltwo}$.





Для любого измеримого множества $X\subset\Bbb R$ введем 


полуторалинейную форму (``скалярное произведение'')


$(f,g)_{X}=\int\limits_{X}f(x)\overline{g(x)}dx$


для функций $f,g\in\Ltwo(\Bbb R)$,


а также полунорму $\norm{f}_{X}=\sqrt{(f,f)_{X}}$.





Пусть $I$, $I_{y}$ --- тождественные операторы в пространства


$\Ltwo(\Bbb R)$ и $\Ltwo(0,y)$ соответственно.








\begin{defnnonum}


Пусть $X\subset\Bbb R$, $f:X\to\Bbb R$. Будем говорить, что функция $f$ 


имеет в точке $x_{0}$ нуль порядка $\nu>0$, если существуют такая 


окрестность $U$ точки $x_{0}$ и такие положительные числа $C_{1}$


и $C_{2}$, что для любого $x\in U-\{x_0\}$ выполняются неравенства


$C_{1}\leqslant\frac{\abs{f(x)}}{\abs{x-x_{0}}^{\nu}}\leqslant C_{2}$.


\end{defnnonum}








\section{Основные результаты} \label{sec:mainres}





Пусть всюду далее $k\in\Lone(\Bbb R)$ --- ядро оператора $A_{y}$, 


определенного равенством \eqref{operatorAy} и действующего в 


пространстве $\Ltwo(0,y)$, $K=\sqrt{2\pi}\cal F k$ --- его символ,


$K$ --- вещественнозначная функция, $\max\limits_{x\in\Bbb R}K(x)=1$.


Пусть, кроме того, $N(y)$ --- 


число положительных собственных чисел оператора $A_{y}$ 


с учетом их кратности; $\lambda_{n}(y)$ ($n\in\Bbb N$, $n\leqslant N(y)$) --- 


упорядоченный набор этих чисел, занумерованных в порядке убывания.


Введем также оператор


$[A f](t)=\int\limits_{\Bbb R}k(t-s)f(s)ds$


свертки $A$, действующий в пространстве $\Ltwo(\Bbb R)$.





В этом разделе сформулируем основные результаты работы, 


доказательство которых будет дано в 


разделах \ref{sec:3}--\ref{sec:mainproof}.








\begin{thms} \label{th:main}


Пусть функция $1-K(x)$ имеет на вещественной прямой единственный нуль 


в точке $x=0$ и его порядок равен $\nu>0$. Тогда для любого $j\in\Bbb N$


существуют такие положительные константы $y_{0}$, $C_{-}$ и $C_{+}$, 


что при всех $y>y_{0}$ выполняется неравенство $N(y)\geqslant j$ и 


для собственного числа $\lambda_{j}(y)$ имеют место оценки


$$


1-\frac{C_{-}}{y^{\nu}} \leqslant \lambda_{j}(y)


\leqslant 1-\frac{C_{+}}{y^{\nu}}.


$$


\end{thms}








\begin{thms} \label{th:general}


Пусть функция $1-K(x)$ имеет на вещественной прямой нули в точках


$x_{l}$, $l\in\set{1,2,\dotsc,n}$, и только в них\/$;$ порядки этих нулей 


равны $\nu_{l}>0$. Обозначим


$\nu_{\max}=\max\limits_{l\in\set{1,2,\dotsc,n}}\nu_{l}$.


Тогда для любого $j\in\Bbb N$ существуют такие положительные константы


$y_{0}$, $C_{-}$ и $C_{+}$, что при всех $y>y_{0}$ выполняется 


неравенство $N(y)\geqslant j$ и для собственного числа $\lambda_{j}(y)$ 


имеют место оценки


$$


1-\frac{C_{-}}{y^{\nu_{\max}}} \leqslant \lambda_{j}(y)


\leqslant 1-\frac{C_{+}}{y^{\nu_{\max}}}.


$$


\end{thms}





В разделах \ref{sec:3} и \ref{sec:4} будем считать, что символ $K$ 


удовлетворяет условию теоремы \ref{th:main}.








\section{Оценка снизу} \label{sec:3}








\begin{lems} \label{le:31}


Существуют такие константы $y_{0}>0$ и $\tld{C}_{-}>0$, что при 


всех $y>y_{0}$ выполняется неравенство $N(y)\geqslant 1$ и для старшего 


собственного числа $\lambda_{1}(y)$ оператора $A_{y}$ имеет место оценка


$$


1-\frac{\tld{C}_{-}}{y^{\nu}}\leqslant\lambda_{1}(y).


$$


\end{lems}








\begin{pf}


Пусть $C>0$ --- такое число, что для любого $x\in\Bbb R$ выполняется 


неравенство $1-K(x)\leqslant C\abs{x}^{\nu}$. Пусть


$m\in\Bbb N$, $m\geqslant\frac{\nu+1}{2}$, $\alpha=\frac{y}{2 m}$. 


Введем функции $\Phi_{m}(x)=C_{m} \big(\frac{\sin x}{x} \big)^{m}$, 


где $C_{m}>0$ --- нормировочная константа, т.\,е. $C_{m}$ --- такое 


число, что $\norm{\Phi_{m}}=1$. Пусть $\phi_{m}=\cal F^{-1}\Phi_{m}$,


$\psi_{m}=\frac{1}{\sqrt{\alpha}}\phi_{m} \big(\frac{x}{\alpha}-m \big)$,


$\Psi_{m}=\cal F\psi_{m}$.





Заметим, что $2\phi_{1}$ --- характеристическая функция сегмента $[-1,1]$, 


поэтому $\supp\phi_{m}\subset[-m,m]$, $\supp \psi_{m}\subset[0,y]$. 


В силу свойств преобразования Фурье $\Psi_{m}(x)=\sqrt{\alpha}


\bold e^{-\bold i \alpha m x}\Phi_{m}(\alpha x)$,


$\norm{\Psi_{m}}=\norm{\psi_{m}}=\norm{\phi_{m}}=\norm{\Phi_{m}}=1$.





Далее имеем


\begin{multline*}


\big((I_{y}-A_{y}) (\psi_{m}|_{[0,y]} ), \psi_{m}|_{[0,y]}\big)_{[0,y]}=


\big([(I-A)\psi_{m}]|_{[0,y]}, \psi_{m}|_{[0,y]} \big)_{[0,y]}=\\


%


=\big((I-A)\psi_{m}, \psi_{m} \big)_{\Bbb R}=


\int_{\Bbb R} (1-K(x))\abs{\Psi_{m}(x)}^{2}dx\leqslant


\alpha C \int_{\Bbb R}\abs{x}^{\nu}\abs{\Phi_{m}(\alpha x)}^{2}\,dx=\\


%


=\frac{C}{\alpha^{\nu}}


\int_{\Bbb R}\abs{x}^{\nu}\abs{\Phi_{m}(x)}^{2}dx=


\frac{C 2^{\nu}m^{\nu}C_{m}^{2}}{y^{\nu}}


\int_{\Bbb R}\abs{x}^{\nu}\abs{\frac{\sin x}{x}}^{2 m}dx=


\frac{\tld{C}_{-}}{y^{\nu}},


\end{multline*}


где


$\tld{C}_{-}=C 2^{\nu}m^{\nu}C_{m}^{2}


\int\limits_{\Bbb R}\abs{x}^{\nu}\big|\frac{\sin x}{x}\big|^{2 m}dx$. 





Отсюда следует, что оператор $I_{y}-A_{y}$ имеет собственное число на отрезке 


$[0,\frac{\Tilde{C}_{-}}{y^{\nu}}]$, а оператор $A_{y}$ --- на отрезке 


$[1-\frac{\Tilde{C}_{-}}{y^{\nu}},1]$. 


Выберем $y_{0}$ так, чтобы $1-\frac{\Tilde{C}_{-}}{y_{0}^{\nu}}>0$. 


Тогда при всех $y>y_{0}$ оператор $A_{y}$ будет иметь положительное 


собственное значение, для которого будет 


выполняться указанная оценка.


\end{pf}








\begin{lems} \label{le:32}


Для любого $j\in\Bbb N$ существуют такие константы $y_{0}>0$ и $C_{-}>0$, 


что при всех $y>y_{0}$ выполняется неравенство $N(y)\geqslant j$ и для 


собственного числа $\lambda_{j}(y)$ оператора $A_{y}$ имеет место оценка


$$


1-\frac{C_{-}}{y^{\nu}}\leqslant\lambda_{j}(y).


$$


\end{lems}








\begin{pf}


Пусть $C$, $m$, $\Phi_{m}$, $C_{m}$, $\phi_{m}$, $\tld{C}_{-}$ --- те же, 


что и при доказательстве леммы \ref{le:31}, $\alpha=\dfrac{y}{2 m j}$.


Введем функции


$\psi_{m,l}=\frac{1}{\sqrt{\alpha}}\phi_{m} \big(\frac{x}{\alpha}


-(2l-1)m \big)$,


$\Psi_{m,l}=\cal F\psi_{m,l}$ для $l\in\set{1,2,\dotsc,j}$.


Тогда $\supp\psi_{m,l}\subset [\frac{l-1}{j}y,\dfrac{l}{j}y ]$,


$\Psi_{m,l}(x)=\sqrt{\alpha}\bold e^{-\bold i


\alpha (2l-1)m x}\Phi_{m}(\alpha x)$, $\norm{\Psi_{m,l}}=


\norm{\psi_{m,l}}=\norm{\phi_{m}}=\norm{\Phi_{m}}=1$.


Заметим, что носители функций $\psi_{m,l}$ при разных $l$ имеют не более 


одной общей точки, а значит, эти функции образуют ортонормированный базис


в $j$-мерном подпространстве пространства $\Ltwo(0,y)$.





Пусть $\gamma_{1},\gamma_{2},\dotsc,\gamma_{j}\in\Bbb C$,


$\psi=\sum\limits_{l=1}^{j}\gamma_{l}\psi_{m,l}$,


$\norm{\psi}=1$, $\Psi=\cal F\psi$.


Действуя так же, как и в лемме \ref{le:31}, получаем


\begin{multline*}


\big((I_{y}-A_{y}) (\psi|_{[0,y]}),


\psi|_{[0,y]} \big)_{[0,y]}=


\int_{\Bbb R} (1-K(x))\abs{\Psi(x)}^{2}\,dx\leqslant\\


%


\leqslant C\int_{\Bbb R}\abs{x}^{\nu}


\bigg|\sum_{l=1}^{j}\gamma_{l}\Psi_{m,l}(x)\bigg|^{2}dx=


C\sum_{l=1}^{j}\abs{\gamma_{l}}^{2} \int_{\Bbb R}\abs{x}^{\nu}


\abs{\Psi_{m,l}(x)}^{2}dx=\\


%


=\alpha C \int_{\Bbb R}\abs{x}^{\nu}\abs{\Phi_{m}(\alpha x)}^{2}\,dx=


\frac{C 2^{\nu} m^{\nu} j^{\nu} C_{m}^{2}}{y^{\nu}}


\int_{\Bbb R}\abs{x}^{\nu}\bigg|\frac{\sin x}{x}\bigg|^{2 m}dx=


\frac{j^{\nu}\tld{C}_{-}}{y^{\nu}}=\frac{C_{-}}{y^{\nu}},


\end{multline*}


где


$C_{-}=j^{\nu}\tld{C}_{-}$.





Таким образом, существует такое подпространство пространства $\Ltwo(0,y)$ 


размерности $j$, что для любой нормированной функции $\psi$ имеем


$((I_{y}-A_{y}) \psi,\psi ) \leqslant\frac{C_{-}}{y^{\nu}}$. 


Из этого следует, что оператор $I_{y}-A_{y}$ имеет на отрезке 


$ [0,\frac{C_{-}}{y^{\nu}} ]$ 


не менее $j$ собственных чисел, а оператор $A_{y}$ --- не менее $j$ 


собственных чисел на отрезке $ [1-\frac{C_{-}}{y^{\nu}},1 ]$. 


Выберем $y_{0}$ так, чтобы $1-\frac{C_{-}}{y_{0}^{\nu}}>0$. 


Тогда при всех $y>y_{0}$ оператор $A_{y}$ будет иметь не менее $j$ 


положительных собственнох значений и для $\lambda_{j}(y)$ 


будет выполняться указанная оценка.


\end{pf}








\section{Оценка сверху} \label{sec:4}





Здесь сначала докажем несколько требуемых в дальнейшем вспомогательных 


положений, представляющих, возможно, и некоторый самостоятельный интерес.








\begin{thms} \label{th:41}


Пусть $f\in\Ltwo(\Bbb R)$, $m=\mu(\supp f)<+\infty,\ F=\cal F f$. Тогда 


имеет место оценка


$$


\sup\limits_{x\in\Bbb R}\abs{F(x)}\leqslant


\sqrt{\frac{m}{2\pi}}\norm{F}_{\Ltwo}.


$$


\end{thms}








{\bf Доказательство.} Пусть $x\in\Bbb R$, тогда


$$


\abs{F(x)}=


\frac{1}{\sqrt{2\pi}}\bigg|\int_{\Bbb R}f(t)\bold e^{-\bold i xt}dt\bigg|


\leqslant


\frac{1}{\sqrt{2\pi}}\int_{\supp f}\abs{f(t)}\abs{\bold e^{-\bold i xt}}dt.


$$


Применим неравенство Коши--Буняковского


$$


\abs{F(x)}\leqslant


\frac{1}{\sqrt{2\pi}}\sqrt{\mu(\supp f)}


\sqrt{\int_{\supp f}\abs{f(t)}^{2}\,dt}=


\sqrt{\frac{m}{2\pi}}\norm{f}_{\Ltwo}=


\sqrt{\frac{m}{2\pi}}\norm{F}_{\Ltwo}.\quad\square


$$








\begin{thms} \label{th:42}


Пусть $X\subset\Bbb R$, $X$ измеримо,


$0<\mu(X)<+\infty$, $m=\frac{\pi}{\mu(X)}$.


Тогда для любой такой функции $F\in\Ltwo(\Bbb R)$, что


$\mu (\supp \cal F^{-1}F )\leqslant m$, имеет место неравенство


$$


\frac{1}{\sqrt{2}}\norm{F}_{\Ltwo} \leqslant


\norm{F}_{\Bbb R\setminus X}.


$$


\end{thms}








{\bf Доказательтство.}


Воспользовавшись теоремой \ref{th:41}, получаем


\begin{gather*}


\norm{F}_{\Ltwo}^{2}=


\norm{F}_{\Bbb R\setminus X}^{2}+\int_{X}\abs{F(x)}^{2}\,dx\leqslant


\norm{F}_{\Bbb R\setminus X}^{2}+\mu(X)\frac{m}{2\pi}\norm{F}_{\Ltwo}^{2},\\


%


\bigg(1-\mu(X)\frac{m}{2\pi} \bigg)\norm{F}_{\Ltwo}^{2}\leqslant


\norm{F}_{\Bbb R\setminus X}^{2},\quad


\frac{1}{2}\norm{F}_{\Ltwo}^{2}\leqslant \norm{F}_{\Bbb R\setminus X}^{2}.	


\quad\square


\end{gather*}








\begin{lems} \label{le:41}


Для любого $y_{0}>0$ существует такая константа $C_{+}>0$, что при 


всех $y>y_{0}$ для любого собственного числа $\lambda$ оператора $A_{y}$ 


имеет место оценка


$$


\lambda\leqslant 1-\dfrac{C_{+}}{y^{\nu}}.


$$


\end{lems}








\begin{pf}


Выберем константу $C>0$ так, чтобы для всех $x\in\Bbb R$ выполнялось 


неравенство


$$


1-K(x)\geqslant C\min\bigg\{\abs{x}^{\nu},


\bigg(\frac{\pi}{2y_{0}} \bigg)^{\nu}\bigg\}.


$$


Пусть $y>y_{0}$, $f\in\Ltwo(\Bbb R)$, $\supp f\subset[0,y]$,


$\norm{f}=1$, $F=\cal F f$,


$\delta=\frac{\pi}{2 y}$, $X=[-\delta,\delta]$. Воспользовавшись 


теоремой \ref{th:42}, получим оценку


\begin{multline*}


\big((I_{y}-A_{y}) (f|_{[0,y]} ), f|_{[0,y]} \big)_{[0,y]}=


\int_{\Bbb R}(1-K(x))\abs{F(x)}^{2}dx\geqslant\\


%


\geqslant C\delta^{\nu}\int_{\abs{x}>\delta}\abs{F(x)}^{2}dx=


C\delta^{\nu}\norm{F}_{\Bbb R\setminus X}^{2}\geqslant


C \bigg(\frac{\pi}{2 y} \bigg)^{\nu}\frac{1}{2}\norm{F}_{\Ltwo}^{2}\geqslant


\frac{C_{+}}{y^{\nu}},


\end{multline*}


где $C_{+}=\frac{C\pi^{\nu}}{2^{\nu+1}}$.


Следовательно, старшее собственное число оператора $A_{y}$ не 


превышает $1-\dfrac{C_{+}}{y^{\nu}}$, но тогда эта оценка верна и для 


всех остальных собственных чисел.


\end{pf}








\section{Доказательство теорем \ref{th:main} и \ref{th:general}} \label{sec:mainproof}





Теорема \ref{th:main} непосредственно следует из 


лемм \ref{le:32} и \ref{le:41}.





Доказательство теоремы \ref{th:general} отличается от доказательства 


теоремы \ref{th:main} лишь техническими деталями. Для доказательства 


леммы, аналогичной лемме \ref{le:32}, достаточно выбрать любую точку, в которой 


функция $1-K(x)$ имеет нуль порядка $\nu_{\max}$, и сдвинуть функции


$\Phi_{m}$ так, чтобы они достигали максимума в этой точке. Для 


получения оценки сверху (аналога леммы \ref{le:41}) следует в качестве 


множества $X$ взять объединение достаточно малых окрестностей всех 


нулей функции $1-K(x)$.








\begin{thebibliography}{9}





\bibitem{li:Serra}


Serra S.


{\em On the extreme eigenvalues of Hermitian $($block\/$)$ Toeplitz matrices}


// Linear Algebra and its Applications. -- 1997. -- V.\,270. -- P.\,109--129.


\bibitem{li:Szego1920}


Szeg\"o G.


{\em Beitr\"age zur Theorie der Toeplitzschen Formen}, I


// Math. Ztschr. -- 1920. -- \No\,6. -- P.\,167--202.


\bibitem{li:Szego1921}


Szeg\"o G.


{\em Beitr\"age zur Theorie der Toeplitzschen Formen}, II


// Math. Ztschr. -- 1921. -- \No\,9. -- P.\,167--190.


\bibitem{li:Szego1952}


Szeg\"o G.


{\em On certain Hermitian forms associated with the Fourier series


of a positive function}


// Festskrift Marcel Riesz. -- Lund. -- 1952. -- P.\,228--238.


\bibitem{li:SzegoForms}


Grenander U., Szeg\"o G.


{\em Toeplitz forms and their applications}. --


New~York: Chelsea, 1984. -- 245~p.


\bibitem{li:SzegoRus}


Гренандер У., Сег\"е Г.


{\em Теплицевы формы и их приложения}. --


М.: Ин. лит., 1961. -- 308~с.





\end{thebibliography}





\vskip 2cm





\post{\it Ростовский государственный}{\it Поступила}


\post{\it университет}{16.06.2003}





\label{end}


\end{document}


